\subsection{Schnittpunktberechnung mit einem LGS}
\label{subsec:analytic_geometry_line_plane_intersection}

Die Anzahl gemeinsamer Punkte entspricht der Anzahl der Lösungen des beim
Gleichsetzen entstehenden LGS. Damit lassen sich alle drei Lagefälle mit einem
einheitlichen Verfahren unterscheiden.

\begin{center}
  \CachedTikzFigure{content/resources/chapters/analytic_geometry/tikz/line_plane_relationships.tex}
\end{center}

\begin{KeyPointBox}{Drei mögliche Lagebeziehungen}{analytic_geometry_line_plane_cases}
  Für eine Gerade $g$ und eine Ebene $E$ gibt es genau drei Fälle:
  \[
    \begin{array}{c@{\qquad\Longleftrightarrow\qquad}c}
      \text{genau eine Lösung} & g\text{ schneidet }E,\\[0.5em]
      \text{keine Lösung} & g\parallel E\text{ ohne gemeinsamen Punkt},\\[0.5em]
      \infty\text{-viele Lösungen} & g\subset E.
    \end{array}
  \]
\end{KeyPointBox}

\begin{TheoremBox}{Gleichsetzen der Parameterdarstellungen}{analytic_geometry_equating_line_plane}
  Für
  \[
    E:\ \vec x=\vec p+\lambda\vec a+\mu\vec b,
    \qquad
    g:\ \vec x=\vec q+\nu\vec c
  \]
  beschreiben die Lösungen von
  \[
    \vec p+\lambda\vec a+\mu\vec b
    =\vec q+\nu\vec c
  \]
  genau die gemeinsamen Punkte von $E$ und $g$.
\end{TheoremBox}

\begin{ExampleBox}{Schnittpunkt von Gerade und Ebene}{analytic_geometry_line_plane_intersection_example}
  Gegeben sind
  \[
    E:\quad
    \vec x=\begin{pmatrix}6\\3\\-5\end{pmatrix}
      +\lambda\begin{pmatrix}1\\1\\-2\end{pmatrix}
      +\mu\begin{pmatrix}-1\\3\\1\end{pmatrix}
  \]
  und
  \[
    g:\quad
    \vec x=\begin{pmatrix}6\\7\\-6\end{pmatrix}
      +\nu\begin{pmatrix}7\\1\\4\end{pmatrix}.
  \]
  Durch Gleichsetzen und Ordnen nach den Parametern folgt
  \begin{ClozeArea}[34]{12}
    \[
      \lambda\begin{pmatrix}1\\1\\-2\end{pmatrix}
      +\mu\begin{pmatrix}-1\\3\\1\end{pmatrix}
      -\nu\begin{pmatrix}7\\1\\4\end{pmatrix}
      =\begin{pmatrix}0\\4\\-1\end{pmatrix}
    \]
    und damit das LGS
    \[
      \left[
        \begin{array}{rrr|r}
          1&-1&-7&0\\
          1& 3&-1&4\\
         -2& 1&-4&-1
        \end{array}
      \right].
    \]
  \end{ClozeArea}
  Die Determinante der Koeffizientenmatrix ist $-66\ne0$. Somit existiert genau
  ein Parametersatz und daher genau ein Schnittpunkt.
\end{ExampleBox}

\clearpage
\begin{tcolorbox}[
  scriptbox,
  breakable=false,
  borderline west={1.2pt}{0pt}{black!40},
  title={Beispiel~\ref{ex:analytic_geometry_line_plane_intersection_example}: Schnittpunkt von Gerade und Ebene -- Fortsetzung}
]
  Das Gauß-Verfahren liefert die Parameter:
  \begin{ClozeArea}[34]{18}
    \[
      \begin{array}{r@{\quad}l@{\qquad}l}
        \phantom{\Longleftrightarrow}
        &\left[
          \begin{array}{rrr|r}
            1&-1&-7&0\\
            1& 3&-1&4\\
           -2& 1&-4&-1
          \end{array}
        \right]
        &\TransformOp{Z_2-Z_1,\;Z_3+2Z_1}\\[1.3em]
        \Longleftrightarrow
        &\left[
          \begin{array}{rrr|r}
            1&-1& -7& 0\\
            0& 4&  6& 4\\
            0&-1&-18&-1
          \end{array}
        \right]
        &\TransformOp{4Z_3+Z_2}\\[1.3em]
        \Longleftrightarrow
        &\left[
          \begin{array}{rrr|r}
            1&-1& -7&0\\
            0& 4&  6&4\\
            0& 0&-66&0
          \end{array}
        \right]
        &\text{Zeilenstufenform}
      \end{array}
    \]
    \[
      -66\nu=0\Longrightarrow\nu=0,\qquad
      4\mu+6\nu=4\Longrightarrow\mu=1,\qquad
      \lambda-\mu-7\nu=0\Longrightarrow\lambda=1.
    \]
  \end{ClozeArea}
  Einsetzen in beide Darstellungen bestätigt denselben Punkt:
  \begin{ClozeArea}[34]{7}
    \[
      \begin{aligned}
        E:\quad
        \vec x
          &=\begin{pmatrix}6\\3\\-5\end{pmatrix}
            +\begin{pmatrix}1\\1\\-2\end{pmatrix}
            +\begin{pmatrix}-1\\3\\1\end{pmatrix}
            =\begin{pmatrix}6\\7\\-6\end{pmatrix},\\[1em]
        g:\quad
        \vec x
          &=\begin{pmatrix}6\\7\\-6\end{pmatrix}
            +0\begin{pmatrix}7\\1\\4\end{pmatrix}
            =\begin{pmatrix}6\\7\\-6\end{pmatrix}.
      \end{aligned}
    \]
  \end{ClozeArea}
  Der Schnittpunkt ist somit
  \[
    S=\Cloze{4cm}{(6,7,-6)}.
  \]
\end{tcolorbox}
